\section[Five consecutive lower-bounded event packets]{Five consecutive lower-bounded event packets}
\label{sec:five-retained-events}

The following argument permits every retained event coefficient to be
unbounded. It requires no canonical omission, actual-contact coordinates,
or source-shape entry. Five consecutive events in the eight-cycle suffice.

\begin{theorem}[Five consecutive lower bounds]
\label{thm:five-retained-events}
Let the original canonical source be physical, with $d>0$ and all eight
event coefficients nonnegative. Suppose
\[
 d/m,\ d/c,\ d/A,\ d/D,\ d/\alpha
 \quad\hbox{are positive bounded coefficient elements}.
\]
Then the entire original parent is tensor SOS at every independent
projective measurement pair and every real or complex projective state.
There is no upper bound on any event coefficient divided by $d$.
By exact CHSH relabeling the same conclusion holds if any five
consecutive event coefficients in the cycle
$(\lambda,D,c,A,\alpha,m,\beta,n)$ have these lower bounds.
\end{theorem}

\begin{proof}
Divide by the positive pure coefficient $d$, so $d=1$.
Let $\mu=\min(1,m,c,A,D,\alpha)$. It is a positive bounded unit.
Select the parent $F_\mu$ with $m=c=A=D=\alpha=\mu$ and
$\lambda=n=\beta=0$. Keep every nonnegative excess and all three
unselected coefficients as their original event squares. This uses only
pure coefficient square roots. It makes no assertion that $F_\mu$ is
globally physical.

At noncommuting residue tests apply
the inherited noncommuting-residue theorem to the full physical source. It
performs its own positive coefficient normalization if the event
weights are unbounded.
It remains to certify $F_\mu$ near every commuting residue test. Its
coefficients are bounded, so a positive-definite real residue Gram is
sufficient by the bounded positive-residue Gram lemma.

Write $\mathcal H$ for the signed core with $d=1$ and
$D_u=1+u^2,D_v=1+v^2$. On the four open commuting axes its exact
homogeneous restrictions are
\[
\begin{aligned}
\mathcal H|_{u=0}&=2D_v(X^2+Y^2),\\
\mathcal H|_{v=0}&=2\{(uW-Y)^2+(X+uZ)^2\},\\
\mathcal H|_{u=\infty}&=2\{(W+vX)^2+(vY-Z)^2\},\\
\mathcal H|_{v=\infty}&=2D_u(X^2+Y^2).
\end{aligned}                                                   \tag{NW1}
\]
The remaining finite angle has nonzero residue. At $u=0$, the retained
$m,D$ packets force $W=Z=0$ on the core kernel. At $v=0$, the retained
$A,\alpha$ packets force $Y=Z=0$, and then $W=X=0$. At $u=\infty$,
the retained $c,A$ packets force $W=Y=0$, and then $X=Z=0$. At
$v=\infty$, the retained $c,D$ packets force $W=Z=0$ after the core
has forced $X=Y=0$. Thus the
selected residue state matrix is positive definite on every open axis,
for either sign of the other finite angle.

Normalize a nonzero testing state by a coordinate of largest absolute
value. Its coordinates are then bounded and one is a unit. At the four
double corners, write $F_{ab}$ for the homogeneous restriction of
$F_\mu$ to $(u,v)=(a,b)$. Directly,
\[
\begin{aligned}
F_{00}&=(2+4\mu)X^2+(2+12\mu)Y^2+4\mu Z^2,\\
F_{\infty0}&=(2+8\mu)(W^2+Z^2)+4\mu Y^2,\\
F_{0\infty}&=4\mu(W^2+Z^2)+2X^2+(2+12\mu)Y^2,\\
F_{\infty\infty}&=4\mu W^2+(2+4\mu)X^2+(2+12\mu)Y^2.
\end{aligned}                                                    \tag{NW2}
\]
Here infinity means the polynomial homogeneous chart, obtained by
dividing the finite parent by the appropriate pure-testing squared
angle. Every form is strict except its explicitly missing state
coordinate. Consequently the only nontrivial residue states are the
$W$ ray at $(0,0)$, the $X$ ray at $(\infty,0)$,
and the $Z$ ray at $(\infty,\infty)$.

At $(0,0)$ normalize $W=1$. The five local coordinates
$(u,v,X,Y,Z)$ are infinitesimal. The complete polynomial $F_\mu$
has no constant or linear term, and its degree-two part is
\[
 (2+8\mu)(u-Y)^2+2X^2
       +4\mu\{(v-X)^2+Y^2+Z^2\}.                    \tag{NW3}
\]
Its residue is positive definite: vanishing successively gives
$X=Y=Z=v=u=0$.

At the corner $(\infty,0)$ normalize $X=1$ and put $p=1/u$.
The five infinitesimal coordinates are $(p,v,W,Y,Z)$. The complete
homogeneous polynomial again has no terms below degree two, and its
quadratic part is
\[
 (2+4\mu)(W+v)^2+(2+4\mu)(p+Z)^2
                         +4\mu(W^2+Y^2+Z^2).        \tag{NW4}
\]
This is positive definite. At the corner $(\infty,\infty)$,
normalize $Z=1$ and put $p=1/u,q=1/v$. On $(p,q,W,X,Y)$ the exact
quadratic part is
\[
 (2+8\mu)(q-Y)^2+2X^2
                         +4\mu\{(p-X)^2+Y^2+W^2\}. \tag{NW5}
\]
It is positive definite as well. One can obtain (NW5) directly from
(NW3) by the exact symmetry
$(u,v;W,X,Y,Z)\mapsto(1/v,1/u;Z,X,Y,W)$, which permutes the five
equal retained weights.

For completeness, these quadratic calculations control the entire
polynomial. In each local chart, allocate every higher monomial to a
product of two of the five coordinate roots. The remaining monomial
has positive degree, hence is infinitesimal; its coefficient is bounded
because it is a real polynomial in the bounded unit $\mu$. Symmetrizing
gives an exact finite five-root Gram whose residue is (NW3), (NW4), or
(NW5). No inversion, root, or comparison mixing the two fields occurs.
The bounded positive-residue Gram lemma therefore certifies the full
selected polynomial, including all higher terms.

Restore the excess event squares, the positive pure-testing homogeneous
and projective state factors, and the positive coefficient $d$. The
zero state is immediate. Splitting a complex state into real and
imaginary parts proves the complex assertion.

Finally the exact full-parent Alice reflection acts on the displayed
cycle indices as $i\mapsto7-i$, and the swap/reflection
$(u,v;W,X,Y,Z)\mapsto(1/v,1/u;Z,X,Y,W)$ acts as $i\mapsto6-i$.
Their composition gives a one-step cycle rotation. Their full-parent
identities have only positive pure-testing homogeneous factors, so they
transport the five-packet conclusion to every cyclic translate.
\end{proof}

This theorem is a sufficient source domain, not an exhaustive zero-$d$
source cover. It requires five selected events to be positive and
noninfinitesimal relative to $d$. Its all-tests conclusion is independent
of actual contact angles and of all original event hierarchies. In
particular, it contains both C1 and C2 sources with all six retained
events bounded below relative to $d$.
